\ficha{Elevação da taxa e do limite, 1910}

\bloco{Fontes lidas}

\textit{Documentos Parlamentares}, volume de 1910. Li na íntegra as duas
exposições do ministro Leopoldo de Bulhões que acompanham as mensagens de Nilo
Peçanha, de 22 de abril (p.~9-15) e de 8 de novembro (p.~23-33); o parecer de
Barbosa Lima e o substitutivo de Galeão Carvalhal, de 28 e 29 de abril
(p.~15-22); o discurso de Carvalhal de 22 de novembro (p.~47-59); o
requerimento de Alcindo Guanabara de 23 de novembro (p.~59-60); os discursos de
Pandiá Calógeras de 7 e 8 de dezembro (p.~186-226) e de Cincinato Braga de 14 de
dezembro (p.~247-321); e a carta de Joaquim Murtinho lida por Azeredo no Senado
em 27 de dezembro (p.~401-409). Da imprensa, o \textit{Correio Paulistano} e o
\textit{Correio da Manhã} de 30 de abril de 1910.

\bloco{Síntese}

Em maio de 1910 os depósitos atingem o limite de 20 milhões de libras e as
emissões cessam. O objeto do desacordo deixa de ser o padrão e passa a ser o
número da taxa e o tamanho do limite. De um lado ficam os que querem manter 15
e ampliar o limite, com a bancada paulista à frente; de outro, os que querem a
alta, para 18 ou para a extinção do aparelho. A taxa de 16 passa como
transação. O substitutivo da Comissão de Finanças é aprovado na Câmara em 18 de
dezembro por 135 a 10 (p.~383): taxa de 16, limite de 60 milhões e restauração
dos fundos de 1899.

\bloco{O governo: de 16 a 18}

A exposição de abril propõe elevar a taxa a 16, suprimir o limite, delegar ao
Executivo as elevações seguintes e devolver o fundo de garantia à função de
1899. O juízo do ministro sobre a lei de 1906 é severo:

\cit[P:2:19]{foi, no caso, uma violencia exercida sobre a valorização monetaria
do papel}{\dpb{p.~13}}

A Caixa é aceita como emissor contra depósito integral, e a taxa só importa ao
papel: \tr[P:2:20]{o cambio não interessa ao ouro, só nos interessa a nós}
\dpb{p.~14}. Em novembro, na última semana do governo, Bulhões emenda a
proposta, trocando a taxa de 16 \tr[P:2:39]{pela de 18 d., actualmente preferida
pela nossa situação economica} \dpb{p.~33}. Identifica os que pedem taxa menor
com os detentores de café que retêm as letras de exportação:

\cit[P:2:34]{São elles, os baixistas - que falsificam a situação
economica}{\dpb{p.~28}}

E defende a retirada de um milhão de libras da Caixa pelo Banco do Brasil,
contra os que tinham o depósito por intocável:

\cit[P:2:36]{ha entre nós fetichistas do metal amarello, que imaginam intangível
o deposito da Caixa}{\dpb{p.~30}}

\bloco{Os opositores de 1906 diante da reforma}

Barbosa Lima, relator, declara-se \tr[P:2:27]{adversario irreductivel daquelle
artificio monetario} \dpb{p.~21} e não toma a iniciativa de propor a nova taxa.
A elevação a 16 entra no projeto por aditivo de Homero Baptista, votado na
Comissão por Barbosa Lima e Paula Ramos, com o único voto contrário de Galeão
Carvalhal (\jor{cp19100430}{Correio Paulistano}{30 abr. 1910}{2}). O relator
recusa ampliar o limite mantendo 15, o que seria a
quebra definitiva do padrão: \tr[P:2:22]{Esse ideal dos valorizadores á rebours
não lograria transformar-se em lei} \dpb{p.~16}. Affonso Costa pede 18 e se
apresenta como \tr[P:2:98]{Inimigo que fui da creação da Caixa} \dpb{p.~92}.
Alcindo Guanabara, que em 1906 queria 12, defende agora a taxa vigente:

\cit[P:2:66]{a unica taxa possivel é a de 15 dinheiros por 1\$000}{\dpb{p.~60}}

\bloco{Os defensores dos 15}

Galeão Carvalhal fala pela bancada paulista e liga a posição ao Convênio:

\cit[P:2:54]{Nós, os paulistas, pugnamos até este momento pela execução do
programma, que foi adoptado desde 1906 com a celebração do Convenio de
Taubaté}{\dpb{p.~48}}

Recusa o rótulo de baixista, \tr[P:2:58]{não estamos pleiteando o cambio baixo;
defendemos a estabilidade na taxa de 15} \dpb{p.~52}, e desconfia do fundo que
o governo quer restaurar: \tr[P:2:60]{Não acreditei nas maravilhas do fundo de
garantia} \dpb{p.~54}. Atribui a David Campista a mesma posição em 1910: ele
\tr[P:2:57]{foi o primeiro a reconhecer a conveniencia da permanencia da taxa de
15 dinheiros} \dpb{p.~51}. Rui Barbosa, na plataforma civilista citada pelo
\textit{Correio da Manhã}, vota no mesmo sentido: \tr[I:per089842_1910_03208:1]{Este
o meu voto, e a elle junto o de que não se altere a taxa de 15}
(\jor{cm19100430}{Correio da Manhã}{30 abr. 1910}{1}).

Cincinato Braga dá ao argumento a forma mais desenvolvida. Não quer a baixa:
\tr[P:2:305]{O que queremos é cambio estavel} \dpb{p.~299}. Quer a alta, mas
gradual e amarrada ao crescimento do fundo de garantia: \tr[P:2:304]{Quero-a
conquistada degráo por degráo} \dpb{p.~298}. E descreve a valorização rápida
como transferência de renda do campo para a cidade:

\cit[P:2:296]{defraudação immoral das classes camponezas}{\dpb{p.~290}}

Na Comissão de Finanças, em abril, segundo o relato do \textit{Correio
Paulistano}, ele próprio \tr[I:per090972_1910_16787:2]{diz que foi inimigo da
Caixa, mas lealmente reconhece} os serviços que ela prestou
(\jor{cp19100430}{Correio Paulistano}{30 abr. 1910}{2}).

\bloco{Os defensores da alta}

Pandiá Calógeras propõe a extinção: \tr[P:2:252]{Fica extincta a Caixa de
Conversão} \dpb{p.~246}. Sustenta que os depósitos cresceram por arbitragem,
fator que \tr[P:2:195]{nada tinha que ver com os saldos do Brasil} \dpb{p.~189},
e põe a questão em termos distributivos: entre os que lucram com o câmbio baixo
e os que ganham com a alta, é mais razoável que \tr[P:2:225]{soffram 5 \% do
que os 95 \% restantes} \dpb{p.~219}. Daí a imagem:

\cit[P:2:228]{E' uma nova campanha abolicionista que se enceta}{\dpb{p.~222}}

Josino de Araujo, defensor dos 15, devolve a metáfora: \tr[P:2:375]{os escravos
são, hoje, os fazendeiros} \dpb{p.~369}.

\bloco{Senado: a carta de Murtinho}

Murtinho, doente, manda ler uma carta contra o limite de 60 milhões. Para ele
os que desejam a valorização estão sendo iludidos

\cit[P:2:408]{pelos que por interesse de classe desejam exactamente o inverso
daquillo que pregam, isto é, a desvalorização da moeda}{\dpb{p.~402}}

Azeredo, que lê a carta, quer 17 e aceita 16; diz que os senadores paulistas
aceitam 16 \tr[P:2:411]{uma vez que não podem conseguir a conservação da taxa
de 15} \dpb{p.~405} e que o limite de 60 milhões seria \tr[P:2:412]{quasi que a
quebra do padrão} \dpb{p.~406}.

\bloco{Achados}

\begin{enumerate}[leftmargin=*,nosep]
\item Em 1910 o eixo é 15 contra a alta, com 16 como transação. Os opositores
de 1906 defendem a letra da lei de 1906 contra a ampliação do limite, ou pedem
taxa mais alta. É mudança de objeto, não de lado.
\item Alcindo Guanabara é coerente como estabilizador a taxa baixa: 12 em
1906, 15 em 1910. Na ficha 6 ele volta a dizer que a taxa certa de 1906 era 12.
\item Os dois campos nomeiam interesses de classe. Bulhões fala em detentores
de café, Murtinho em interesse de classe, Calógeras em 95 por cento contra 5;
Cincinato inverte os termos, com classes camponesas contra urbanas. Calógeras e
Josino disputam a imagem da abolição.
\item A posição do governo sobre o depósito é a de que o ouro pertence ao
depositante, que pode retirá-lo. O mesmo princípio serve, em 1914, para
condenar a suspensão do troco (ficha 7).
\item Campista aparece em 1910 a favor da permanência dos 15, mas só pela voz
de Carvalhal.
\end{enumerate}

\bloco{Limites}

Dos discursos de 1910, só os de Carvalhal, Calógeras e Cincinato e as peças do
governo foram lidos por inteiro. A posição de Rui Barbosa vem da plataforma
citada em jornal e por Carvalhal na tribuna. A autodescrição de Cincinato como
ex-inimigo da Caixa não tem correspondência em texto de 1906: ali ele só
aparece votando \textit{sim}.
